% 分野別類題: 完全微分方程式（1階: M dx + N dy = 0）
% 完全性の判定とポテンシャル関数の構成、積分因子による完全化を含む 3 問。
% コンパイル: TEXINPUTS="../../../00_common:$TEXINPUTS" latexmk -xelatex problems.tex
\documentclass[a4paper,11pt]{article}
\usepackage{preamble}

\begin{document}

\kamokumei{類題演習 --- 完全微分方程式}

\shijibun{次の常微分方程式を解き、導出過程を示せ。（$M(x,y)\,dx + N(x,y)\,dy = 0$ が完全微分方程式であるとは、$\dfrac{\partial M}{\partial y} = \dfrac{\partial N}{\partial x}$ を満たし、ある関数 $F(x,y)$ が存在して $dF = M\,dx + N\,dy$ となることをいう。このとき一般解は $F(x,y) = C$ で与えられる。完全でない場合は、適当な積分因子を掛けて完全微分方程式に直してから解く。）}

\shomon{問1.}{次の常微分方程式が完全微分方程式であることを確かめ、一般解を求めよ。}
\[
(2xy + 3x^{2})\,dx + (x^{2} + 2y)\,dy = 0
\]

\shomon{問2.}{次の常微分方程式は完全ではないが、$x$ のみに依存する積分因子 $\mu(x)$ を掛けると完全微分方程式になる。$\mu(x)$ を求め、一般解を求めよ。}
\[
(3xy + y^{2})\,dx + (x^{2} + xy)\,dy = 0
\]

\shomon{問3.}{次の初期値問題を解け。}
\[
(2x + y\cos x)\,dx + (\sin x - 2y)\,dy = 0, \qquad y(0) = 3
\]

\end{document}
